\documentclass[a4paper]{article}

%------------------------------------------------
% الحزم
%------------------------------------------------

\usepackage{enumitem}
\usepackage[margin=2.2cm]{geometry}
\usepackage{amsmath,amssymb}
\usepackage{graphicx}
\usepackage{array}
\usepackage{fancyhdr}
\usepackage[table]{xcolor}
\usepackage{tikz}
\usetikzlibrary{positioning,arrows.meta,shapes.geometric,calc}
\usepackage[skins]{tcolorbox}
\usepackage{polyglossia}

%------------------------------------------------
% اللغات والخطوط
%------------------------------------------------

\setdefaultlanguage[
    calendar=gregorian,
    numerals=maghrib
]{arabic}

\setotherlanguage{english}

% إذا لم تكن الخطوط مثبتة في الجهاز، تُؤخذ من المجلد fonts/
\IfFontExistsTF{Amiri}{%
    \newfontfamily\arabicfont[Script=Arabic, Scale=1.15]{Amiri}
}{%
    \newfontfamily\arabicfont[
        Script=Arabic, Scale=1.15,
        Path=fonts/, Extension=.ttf,
        UprightFont=*-Regular, BoldFont=*-Bold,
        ItalicFont=*-Italic, BoldItalicFont=*-BoldItalic
    ]{Amiri}
}

\IfFontExistsTF{AlBattar}{%
    \newfontfamily\arabicfontsf[Script=Arabic, Scale=1.2]{AlBattar}
    % خط العناوين (يعمل حتى بعد استعمال \textenglish)
    \newfontfamily\khatt[Script=Arabic, Scale=1.2]{AlBattar}
}{%
    \newfontfamily\arabicfontsf[
        Script=Arabic, Scale=1.2, Path=fonts/,
        BoldFont=ae_AlBattar.ttf
    ]{ae_AlBattar.ttf}
    % خط العناوين (يعمل حتى بعد استعمال \textenglish)
    \newfontfamily\khatt[
        Script=Arabic, Scale=1.2, Path=fonts/,
        BoldFont=ae_AlBattar.ttf
    ]{ae_AlBattar.ttf}
}

%------------------------------------------------
% الألوان: لون رئيسي هادئ + لون تنبيه فقط
%------------------------------------------------

\definecolor{asasi}{HTML}{3D6A8E}   % اللون الرئيسي: للرسوم والمخططات
\definecolor{tanbih}{HTML}{C0603A}  % لون التنبيه: للتحذيرات والحالات السلبية
\definecolor{ramadi}{HTML}{5F6B73}  % رمادي للنصوص الثانوية والأسهم
\definecolor{dakin}{HTML}{263238}   % لون العناوين

% أسماء قديمة تُحوَّل إلى اللونين فقط
\colorlet{azraq}{asasi}
\colorlet{fayruzi}{asasi}
\colorlet{banafsaji}{asasi}
\colorlet{akhdar}{asasi}
\colorlet{ahmar}{tanbih}
\colorlet{burtuqali}{asasi}

%------------------------------------------------
% الصناديق
%------------------------------------------------

\tcbset{
    asl/.style={
        enhanced,
        arc=2mm,
        boxrule=0.6pt,
        colback=white,
        colframe=black!18,
        colbacktitle=black!5,
        coltitle=dakin,
        titlerule=0pt,
        left=8pt, right=8pt, top=5pt, bottom=5pt,
        fonttitle=\khatt\bfseries,
        before skip=8pt, after skip=8pt
    }
}

% أهداف الدرس
\newtcolorbox{ahdaf}{asl, title={أهداف الدرس}}

% تعريف (خط رئيسي رفيع على الجانب)
\newtcolorbox{taarif}[1][تعريف]{asl, title={#1},
    borderline east={2pt}{0pt}{asasi}}

% مثال
\newtcolorbox{mithal}[1][مثال]{asl, title={#1}}

% تنبيه (اللون الوحيد البارز في النص)
\newtcolorbox{intabih}[1][انتبه!]{asl, colback=tanbih!5, colframe=tanbih,
    colbacktitle=tanbih, coltitle=white, title={#1}}

% ركن أستاذ اللغة العربية
\newtcolorbox{arabiyya}[1][في قسم اللغة العربية]{asl, title={#1}}

% نص أمر يكتب للأداة
\newtcolorbox{amr}[1][ramadi]{enhanced, arc=1mm, boxrule=0pt,
    colback=black!4, borderline east={2.5pt}{0pt}{#1},
    left=8pt, right=10pt, top=4pt, bottom=4pt,
    before skip=6pt, after skip=6pt}

% الخلاصة
\newtcolorbox{khulasa}{asl, title={خلاصة الدرس}}

% عنوان القسم
\newcounter{qism}
\newcommand{\qism}[1]{%
    \stepcounter{qism}%
    \par\bigskip
    {\large\khatt\bfseries{\color{asasi}\arabic{qism}.}\ {\color{dakin}#1}}\par
    \vspace{-2pt}{\color{black!20}\rule{\linewidth}{0.6pt}}\par\nopagebreak[4]
    \smallskip}

% عنوان فرعي
\newcommand{\far}[1]{\par\smallskip
    {\color{dakin}\khatt\bfseries {\color{asasi}$\blacktriangleleft$}\ #1}\par\nopagebreak[4]}

% كلمة مهمة (خط عريض فقط)
\newcommand{\kalima}[2][]{\textbf{#2}}

%------------------------------------------------
% أنماط الرسوم
%------------------------------------------------

\newcommand{\nodecolor}{}

\tikzset{
    % كتابة النص العربي من اليمين إلى اليسار داخل العقد
    % (لون النص يُعطى بالمفتاح lawn=... وليس text=...)
    every node/.append style={
        execute at begin node={\nodecolor\everypar{\beginR}\beginR},
        execute at end node=\endR},
    lawn/.code={\def\nodecolor{\color{#1}}},
    khatwa/.style={draw=#1, fill=#1!12, very thick, rounded corners=6pt,
        minimum height=1.2cm, text width=2.3cm, align=center,
        font=\khatt\bfseries},
    khatwa/.default=azraq,
    sahm/.style={-{Stealth[length=3mm]}, very thick, color=ramadi},
    raqm/.style={circle, fill=#1, lawn=white, inner sep=2pt,
        font=\small\bfseries},
}

%------------------------------------------------
% إعداد الصفحة
%------------------------------------------------

\renewcommand{\footrulewidth}{1pt}
\renewcommand{\headrulewidth}{1pt}

\rhead{\khatt المدرسة العليا للأساتذة ببوزريعة}
\lhead{\khatt ملحقة بومرداس}

\cfoot{\thepage}
\lfoot{\khatt الأستاذ: د. بن عبد الرزاق يوسف}

\parindent 0pt
\setlength{\headsep}{15pt}
\setlength{\parskip}{4pt}

\pagestyle{fancy}

\setlist[enumerate,1]{label=\arabic*), itemsep=2pt, topsep=3pt}
\setlist[itemize]{itemsep=2pt, topsep=3pt}

% الجداول
\arrayrulecolor{black!22}
\renewcommand{\arraystretch}{1.35}

% رأس جدول ملون
\newcommand{\ras}[1]{\cellcolor{black!7}\color{dakin}\bfseries #1}

% عنوان الدرس
\newcommand{\unwan}[3]{%
    \begin{center}
        {\color{asasi}\rule{\linewidth}{1.2pt}}\par\medskip
        {\khatt\color{ramadi} #1}\par\smallskip
        {\Large\khatt\bfseries\color{dakin} #2}\par\smallskip
        {\khatt\color{ramadi} #3}\par\medskip
        {\color{asasi}\rule{\linewidth}{1.2pt}}
    \end{center}}

%================================================
% بداية الوثيقة
%================================================

\begin{document}

%------------------------------------------------
% العنوان
%------------------------------------------------

\unwan{درس الفصل الخامس}
{البحث عن المعلومات، تحليلها وتلخيصها}
{من السؤال إلى المصدر الموثوق}

\begin{ahdaf}
في نهاية هذا الدرس، تستطيع أن:
\begin{itemize}
    \item تفرّق بين موضوع البحث وسؤال البحث.
    \item تختار كلمات مفتاحية مناسبة.
    \item تميّز بين أنواع المصادر وتقيّم جودتها.
    \item تلخّص المعلومات بطريقة منظمة وأمينة.
\end{itemize}
\end{ahdaf}

%================================================
\qism{مراحل البحث عن المعلومة}
%================================================

البحث الجيد لا يبدأ بمحرك البحث، بل يبدأ بسؤال واضح.
نمرّ بأربع مراحل:

\begin{center}
\begin{tikzpicture}[node distance=7mm]
    \node[khatwa=azraq]  (a) {حدد\\الحاجة};
    \node[khatwa=fayruzi, left=of a]   (b) {ابحث};
    \node[khatwa=burtuqali, left=of b] (c) {قيّم\\المصادر};
    \node[khatwa=akhdar, left=of c]    (d) {لخّص};
    \draw[sahm] (a)--(b); \draw[sahm] (b)--(c); \draw[sahm] (c)--(d);
\end{tikzpicture}
\end{center}

%================================================
\qism{من الموضوع إلى سؤال البحث}
%================================================

\begin{taarif}
\textbf{موضوع البحث:} مجال عام وواسع، مثل: «الذكاء الاصطناعي والتعليم».

\textbf{سؤال البحث:} سؤال محدد نريد الإجابة عنه،
مثل: «كيف يساعد الذكاء الاصطناعي التلميذ على تحسين التعبير الكتابي؟».
\end{taarif}

الموضوع العام يعطي نتائج كثيرة جدًا ومتفرقة.
أما السؤال المحدد فيوجّه البحث نحو معلومات مفيدة.

\begin{center}
\begin{tikzpicture}
    \fill[azraq!15]  (-6,0)   -- (6,0)   -- (4.5,-1.1) -- (-4.5,-1.1) -- cycle;
    \fill[azraq!35]  (-4.5,-1.1) -- (4.5,-1.1) -- (3,-2.2) -- (-3,-2.2) -- cycle;
    \fill[azraq!60]  (-3,-2.2) -- (3,-2.2) -- (1.5,-3.3) -- (-1.5,-3.3) -- cycle;
    \node[font=\khatt\bfseries] at (0,-0.55) {الذكاء الاصطناعي};
    \node[font=\khatt\bfseries] at (0,-1.65) {الذكاء الاصطناعي في التعليم الجامعي};
    \node[font=\khatt\bfseries, lawn=white] at (0,-2.75) {فوائده ومخاطره};
    \node[font=\small, lawn=ramadi, anchor=west] at (6.3,-0.55) {عام جدًا};
    \node[font=\small, lawn=ramadi, anchor=west] at (4.8,-1.65) {أدق};
    \node[font=\small, lawn=azraq, anchor=west] at (3.3,-2.75) {\textbf{دقيق}};
\end{tikzpicture}
\end{center}

\begin{arabiyya}[مثال من اختصاصك]
\textbf{موضوع:} تعليم النحو.\\
\textbf{سؤال 1:} ما الصعوبات التي يواجهها تلاميذ المتوسط في الإعراب؟\\
\textbf{سؤال 2:} هل تساعد الأمثلة من الحياة اليومية على فهم القاعدة؟\\
\textbf{سؤال 3:} كيف يمكن استعمال الذكاء الاصطناعي في تمارين الإعراب؟
\end{arabiyya}

%================================================
\qism{الكلمات المفتاحية}
%================================================

\begin{taarif}
\textbf{الكلمات المفتاحية:} الكلمات الأساسية في سؤال البحث،
نكتبها في محرك البحث للوصول إلى المعلومة.
\end{taarif}

\far{كيف نختارها؟}
\begin{enumerate}
    \item نقرأ سؤال البحث، ونحذف الكلمات العامة مثل: «كيف»، «ما»، «هل».
    \item نحتفظ بالكلمات التي تحمل المعنى.
    \item نضيف المرادفات، والكلمات الإنجليزية المقابلة.
\end{enumerate}

\begin{mithal}
\textbf{السؤال:} كيف يساعد الذكاء الاصطناعي على تحسين التعبير الكتابي؟

\textbf{الكلمات المفتاحية:}
\kalima{الذكاء الاصطناعي} ــ \kalima{التعبير الكتابي} ــ \kalima{الإنشاء} ــ
\textenglish{AI} ــ \textenglish{Writing skills}
\end{mithal}

%================================================
\qism{أنواع المصادر ودرجة الثقة}
%================================================

ليست كل المصادر في درجة الثقة نفسها.
نرتبها من الأكثر موثوقية إلى الأقل:

\begin{center}
\begin{tikzpicture}[
    masdar/.style={draw=#1, fill=#1!15, thick, rounded corners=4pt,
        minimum height=0.85cm, text width=7cm, align=center}]
    \node[masdar=akhdar]    at (0,0)    {مصدر أكاديمي: كتاب محقق، مقال علمي محكّم};
    \node[masdar=fayruzi]   at (0,-1)   {موقع مؤسسة رسمية: وزارة، جامعة};
    \node[masdar=ramadi] at (0,-2)   {مقال صحفي};
    \node[masdar=ahmar!70!burtuqali] at (0,-3) {مدونة أو موقع شخصي};
    \node[masdar=ahmar]     at (0,-4)   {محتوى أنتجه الذكاء الاصطناعي};
    \draw[{Stealth[length=4mm]}-, line width=3pt, akhdar!70!ahmar]
        (4.6,0.3) -- (4.6,-4.3);
    \node[font=\khatt\bfseries, lawn=akhdar, anchor=west] at (4.9,0) {ثقة أكبر};
    \node[font=\khatt\bfseries, lawn=ahmar,  anchor=west] at (4.9,-4) {ثقة أقل};
\end{tikzpicture}
\end{center}

\begin{intabih}
محتوى الذكاء الاصطناعي ليس مصدرًا علميًا.
نستعمله للفهم أو لاقتراح أفكار، ثم نبحث عن المعلومة في مصدر حقيقي.
\end{intabih}

\begin{intabih}[لا تكتفِ بالنتيجة الأولى]
النتيجة الأولى في محرك البحث ليست بالضرورة الأفضل.
قد تكون إعلانًا، أو موقعًا مشهورًا لكنه غير دقيق.
\end{intabih}

%================================================
\qism{معايير تقييم المصدر}
%================================================

قبل أن نعتمد على مصدر، نطرح خمسة أسئلة:

\begin{center}
\rowcolors{2}{black!3}{white}
\begin{tabular}{>{\centering\arraybackslash}m{2.6cm}
                >{\raggedleft\arraybackslash}m{11.3cm}}
\hline
\ras{المعيار} & \ras{السؤال} \\
\hline
\kalima{الكاتب}           & من كتب هذا المحتوى؟ هل هو مختص في المجال؟ \\
\kalima[fayruzi]{الجهة}    & من نشره؟ جامعة، مجلة علمية، أم موقع مجهول؟ \\
\kalima[burtuqali]{التاريخ} & متى نُشر؟ هل المعلومة ما زالت صالحة؟ \\
\kalima[akhdar]{الدقة}     & هل يذكر مراجعه؟ هل نجد المعلومة نفسها في مصدر آخر؟ \\
\kalima[banafsaji]{الهدف}  & هل هدفه الإعلام، أم الإقناع، أم البيع؟ \\
\hline
\end{tabular}
\end{center}

%================================================
\qism{تلخيص المعلومات}
%================================================

\begin{taarif}
\textbf{التلخيص:} إعادة كتابة الأفكار الأساسية لنص ما،
بكلمات أقل وبأسلوبنا الخاص، دون تغيير المعنى.
\end{taarif}

\begin{center}
\begin{tikzpicture}[node distance=6mm]
    \node[khatwa=azraq, text width=2.2cm] (a) {اقرأ النص\\كاملًا};
    \node[khatwa=fayruzi, text width=2.2cm, left=of a] (b) {حدد الأفكار\\الأساسية};
    \node[khatwa=burtuqali, text width=2.2cm, left=of b] (c) {احذف التكرار\\والأمثلة};
    \node[khatwa=banafsaji, text width=2.2cm, left=of c] (d) {اكتب\\بأسلوبك};
    \node[khatwa=akhdar, text width=2.2cm, left=of d] (e) {اذكر\\المصدر};
    \draw[sahm] (a)--(b); \draw[sahm] (b)--(c);
    \draw[sahm] (c)--(d); \draw[sahm] (d)--(e);
\end{tikzpicture}
\end{center}

\begin{arabiyya}[الذكاء الاصطناعي والتلخيص]
يمكن أن نطلب من الأداة تلخيص نص طويل.
لكن يجب أن نقارن الملخص بالنص الأصلي:
هل حذفت فكرة مهمة؟ هل أضافت معلومة غير موجودة؟
\end{arabiyya}

\begin{khulasa}
\begin{itemize}
    \item نبدأ البحث بــ\textbf{سؤال محدد}، لا بموضوع عام.
    \item نختار \textbf{كلمات مفتاحية} دقيقة ونضيف مرادفاتها.
    \item نقيّم المصدر: الكاتب، الجهة، التاريخ، الدقة، الهدف.
    \item نلخّص بأسلوبنا، ونذكر \textbf{المصدر} دائمًا.
\end{itemize}
\end{khulasa}

\end{document}
